\subsection{并集与交集的性质}

如果 \((\mathbf{X}_t)_{t\in \mathbf{I}}\) 并且 \((\mathbf{Y}_t)_{t\in \mathbf{I}}\) 是具有相同指标集的集合族 \(\mathbf{I}\) ，并且如果 \(\mathbf{Y}_t\subset \mathbf{X}_t\) 对于每一个 \(t\in \mathbf{I}\) 那么显然

\[
\bigcup_ {i \in I} Y _ {i} \subset \bigcup_ {i \in I} X _ {i}, \quad \text { 且 } (\text { 若   } I \neq \emptyset) \quad \bigcap_ {i \in I} Y _ {i} \subset \bigcap_ {i \in I} X _ {i}.
\]

让 \((\mathbf{X}_t)_{t\in \mathbf{I}}\) 设为一个集合族。若 \(J\subset I\) ，我们有

\[
\bigcup_ {i \in J} X _ {i} \subset \bigcup_ {i \in I} X _ {i}, \quad \text { 且 } (\text { 若 } J \neq \emptyset) \quad \bigcap_ {i \in J} X _ {i} \supset \bigcap_ {i \in I} X _ {i}.
\]

命题2. 设 \((\mathbf{X}_{\iota})_{\iota \in \mathbf{I}}\) 为一个集合族，其指标集 I 是某个族 \((\mathrm{J}_{\lambda})_{\lambda \in \mathbf{L}}\) 的并集。则

\[
\bigcup_ {i \in I} X _ {i} = \bigcup_ {\lambda \in L} \left(\bigcup_ {i \in J _ {\lambda}} X _ {i}\right),
\]

并且（如果 \(\mathbf{L} \neq \phi\) 并且 \(J_{\lambda} \neq \phi\) 对于每一个 \(\lambda \in \mathbf{L}\) )

\[
\bigcap_ {i \in I} X _ {i} = \bigcap_ {\lambda \in L} \left(\bigcap_ {i \in J _ {\lambda}} X _ {i}\right)
\]

（并集与交集的“结合律”）。

让 \(x\) 为 \(\bigcup_{i \in I} X_i\) . 存在一个指标 \(i \in I\) 使得 \(x \in X_i\) 。由于 \(I\) 是族 \((J_\lambda)_{\lambda \in L}\) ，存在一个指标 \(\lambda \in L\) 使得 \(i \in J_\lambda\) ，由此 \(x \in \bigcup X_i\) ，并且因此

\[
x \in \bigcup_ {\lambda \in L} \left(\bigcup_ {i \in J _ {\lambda}} X _ {i}\right).
\]

反之，设 \(x\) 是此集合的一个元素。存在一个指标 \(\lambda \in L\) 使得 \(x \in \bigcup_{\iota \in J_{\lambda}} X_{\iota}\) ，因此存在一个指标 \(\iota \in J_{\lambda}\) （从而 \(\iota \in I\) ) 使得 \(x \in X_{\iota}\) ；由此得出

\[
x \in \bigcup_ {i \in I} X _ {i}.
\]

现在假设 \(L \neq \emptyset\) 并且 \(J_{\lambda} \neq \emptyset\) 对于每一个 \(\lambda \in L\) 。然后 \(I \neq \emptyset\) 。设 x 为 \(\bigcap_{i \in I} X_{i}\) 。如果 \(\lambda \in L\) ，我们有 \(x \in X_{i}\) 对于每一个 \(i \in J_{\lambda}\) （因为 \(J_{\lambda} \subset I\) ），所以 \(x \in \bigcap_{i \in J_{\lambda}} X_{i}\) 。由于最后的包含关系对所有 \(\lambda \in L\) ，x 属于 \(\bigcap_{\lambda \in L} \left( \bigcap_{i \in J_{\lambda}} X_{i} \right)\) 。反之，设 x 是后一个集合的元素，并设 i 是 I 的任意元素。存在 \(\lambda \in L\) 使得 \(i \in J_{\lambda}\) ；由于 \(x \in \bigcap_{i \in J_{\lambda}} X_{i}\) ，我们有 \(x \in X_{i}\) ，这对所有 \(i \in I\) 。因此 \(x \in \bigcap_{i \in I} X_{i}\) 都成立。这就完成了证明。

对于给定集合的子集族，命题2的第二部分在不对L作限制的情况下仍然成立，且 \(J_{\lambda}\) .

